\documentclass[12pt]{article}
\usepackage[T1]{fontenc}
\usepackage[utf8]{inputenc}
\usepackage{lmodern}
\usepackage{textcomp}
\usepackage{enumitem}
\usepackage{geometry}
\geometry{margin=1in}
\begin{document}
\begin{center}
Answer Key - Sociology - Undergraduate Level Assessment 24 \
\small (Answers are ordered exactly as in the question paper.)
\end{center}

\section*{Multiple Choice}
\begin{enumerate}[label=\arabic*.]
\item \textbf{Answer:} B
\\ \textit{Options:} A) globalization; B) stability and the structure of society; C) secularism and mechanization; D) industrialization
\\ \textit{Note:} Equilibrium theories of social change emphasize stability and the structure of society because they focus on how social systems maintain balance and order, suggesting that changes occur gradually and are often adaptive responses to disruptions rather than abrupt transformations.
\item \textbf{Answer:} D
\\ \textit{Options:} A) rationalization; B) colonization; C) McDonaldization; D) socialization
\\ \textit{Note:} Socialization is the process through which individuals learn and internalize the values, norms, and behaviors necessary to function as members of society, thereby shaping their sense of identity.
\item \textbf{Answer:} D
\\ \textit{Options:} A) the group of structural theories of society that he was reacting against; B) the overall impression of ourselves that we try to give off to others; C) a significant figure in early childhood who teaches us the general values of society; D) an image of how people in the wider society might perceive our behaviour
\\ \textit{Note:} The 'generalized other' in Mead's Symbolic Interactionist theory refers to the internalized societal norms and expectations that shape how individuals perceive themselves and their behavior in relation to the broader community.
\item \textbf{Answer:} C
\\ \textit{Options:} A) they have a deep-rooted sense of ineffectiveness and are striving for control; B) they are in pursuit of the 'perfect' body, fuelled by images of beauty; C) they perceive a part of their body as stigmatizing, in relation to a cultural ideal; D) their male partners pressurize them to look like supermodels
\\ \textit{Note:} The correct answer is supported by the concept that societal standards and cultural ideals create pressures for women to conform to certain beauty norms, leading them to view specific aspects of their bodies as stigmatizing and in need of modification through cosmetic surgery.
\item \textbf{Answer:} A
\\ \textit{Options:} A) socialization into working class families and communities; B) rational calculation of self-interest; C) issue-based concerns about culture and lifestyle; D) false consciousness and the acceptance of hegemonic values
\\ \textit{Note:} The correct answer highlights the key sociological concept of socialization, illustrating that individuals raised in working-class families and communities are likely to adopt the political values and affiliations, such as support for the Labour Party, prevalent within those environments.
\end{enumerate}

\section*{True / False}
\begin{enumerate}[label=\arabic*.]
\setcounter{enumi}{5}
\item \textbf{Answer:} True
\\ \textit{Options:} A) True; B) False
\\ \textit{Note:} The statement is true because sociologists recognize that personal values and cultural contexts shape the interpretation of social phenomena, leading to the understanding that knowledge production is influenced by the subjective perspectives of theorists.
\item \textbf{Answer:} False
\\ \textit{Options:} A) True; B) False
\\ \textit{Note:} The New Labour government, while promoting education reform, favored a market-driven approach that often undermined local education authorities by prioritizing initiatives like academies and private funding over direct support for struggling local authorities.
\item \textbf{Answer:} True
\\ \textit{Options:} A) True; B) False
\\ \textit{Note:} Environmental social movements are considered global because they address issues that transcend national boundaries, such as climate change and biodiversity loss, and they effectively leverage global media platforms to disseminate their messages and mobilize support across diverse populations.
\item \textbf{Answer:} True
\\ \textit{Options:} A) True; B) False
\\ \textit{Note:} Becker's theory of deviance posits that cannabis use is a learned behavior acquired through social interactions and experiences, suggesting that individuals progress through stages of a deviant career where they learn the norms and techniques associated with cannabis use.
\item \textbf{Answer:} True
\\ \textit{Options:} A) True; B) False
\\ \textit{Note:} In modern societies, a person's occupation is often linked to their income, education, and social networks, all of which significantly influence their perceived social status and position within the social hierarchy.
\end{enumerate}

\section*{Long Form Response}
\begin{enumerate}[label=\arabic*.]
\setcounter{enumi}{10}
\item \textbf{Answer:} Max Weber's ideal type of bureaucracy is characterized by several key elements, including a clear hierarchy of authority, a set of formal rules and regulations, a division of labor, impersonal relationships, and merit-based advancement. While extensive paperwork is often associated with bureaucratic organizations, it is not considered a defining characteristic of Weber's model. Instead, the focus is on the structured processes and rational organization that enable efficiency and predictability in administrative functions. Paperwork may be a byproduct of these processes, but it does not inherently define the bureaucratic structure itself. This distinction emphasizes that the essence of bureaucracy lies in its systematic approach rather than the quantity of documentation produced.
\\ \textit{Note:} correct because it accurately identifies the core elements of Weber's ideal type of bureaucracy, emphasizing that its defining characteristics focus on structure and rationality rather than the quantity of paperwork, which can vary across different bureaucratic organizations.
\item \textbf{Answer:} Official statistics on strike action often fail to capture the full scope of labor unrest, primarily because many strikes go unnoticed by employers and the media. This underreporting can be attributed to several factors, including the prevalence of unofficial or spontaneous strikes that do not follow formal procedures for reporting. Additionally, strikes may occur in small, localized settings that lack visibility, making them less likely to attract media attention. Furthermore, employers may underestimate the impact of minor grievances that lead to short-lived strikes, viewing them as insignificant. These limitations in statistical data can obscure the true extent of worker discontent and the dynamics of labor relations, thus hindering a comprehensive understanding of strike action in society.
\\ \textit{Note:} correct because it identifies that many strikes are unofficial or localized, leading to underreporting in official statistics due to a lack of visibility and formal reporting procedures, which ultimately skews the understanding of labor unrest.
\end{enumerate}

\vfill
\noindent\textit{End of Answer Key}
\end{document}